\chapter{Number Formats and Dates}
\label{ch:formats}

\section{Where formats come from}
\index{number formats}

\begin{cccaution}
\textbf{There is no command in \ccprod{} that sets a number format.} Not on a
menu, not on a key. The eight formats described in this chapter are displayed
faithfully when a cell has one, but the only ways for a cell to acquire one are
to be imported from an \fname{.XLSX} file that already had it, or to be loaded
from a \fname{.X16S} workbook that already had it.

Everything you type gets the General format and keeps it.
\end{cccaution}

The same applies to bold, to left and right alignment, and to column widths.
The machinery is all there, the file formats carry it, and the display honours
it --- there is simply no user interface for it yet.

\section{The eight formats}
\index{number formats!list}

\begin{cctablethree}{Format}{Example value}{Displays as}{1.35in}{1.15in}{2.3in}
General & 8.5 & \lit{8.5} \\
General & 1/3 & \lit{.333333333} \\
General & 1200000000 & \lit{1.2E+09} \\
Integer & 3.7 & \lit{4} \\
Decimal, 2 places & 8.5 & \lit{8.50} \\
Currency, 2 places & 8.5 & \lit{\$8.50} \\
Percent, 1 place & 0.125 & \lit{12.5\%} \\
Date & 45000 & \lit{2023-03-15} \\
Time & 0.5 & \lit{.5} \\
Text & 8.5 & \lit{8.5} \\
\end{cctablethree}

\subsection{General}

The default, and what every cell you type gets.

\begin{ccsteps}
\item Zero displays as \lit{0}.
\item Up to nine significant digits, with no trailing zeros. \lit{8.50} typed
into a cell displays as \lit{8.5}.
\item \textbf{No leading zero below 1.} One third is \lit{.333333333} and a
half is \lit{.5}. This is the Commodore convention and it is what the ROM's
number formatter produces.
\item A value of a thousand million or more, or smaller than one hundredth,
switches to scientific notation: \lit{1.2E+15}.
\end{ccsteps}

\subsection{Integer}

Rounded to a whole number and shown with no decimal point. It rounds rather
than truncating, so 3.7 displays as \lit{4}. The stored value is unchanged.

\subsection{Decimal}

A fixed number of decimal places, always shown, up to nine. Unlike General,
this format \emph{does} print a leading zero: 0.5 at two places is \lit{0.50}.

\subsection{Currency}

A dollar sign followed by the value at a fixed number of decimal places.

\begin{ccnote}
\textbf{The currency symbol is always \lit{\$}.} There is no setting for it,
and there are no thousands separators either --- eight thousand four hundred
displays as \lit{\$8400.00}, not \lit{\$8,400.00}. A workbook imported from
Excel with a pounds or euro format is displayed in dollars.
\end{ccnote}

\subsection{Percent}

The value multiplied by a hundred, shown to a fixed number of places, with a
per-cent sign. 0.125 at one place is \lit{12.5\%}.

\subsection{Date}

The value read as a day number and shown as \lit{YYYY-MM-DD}. See below.

\subsection{Time}

\begin{ccnote}
\textbf{The time format is unfinished.} A cell marked as a time displays its raw
number --- 0.5 shows as \lit{.5}, not as \lit{12:00:00}. This is deliberate:
the choice was to show something visibly wrong rather than something plausibly
wrong. If you import a workbook with times in it, the times will look like
fractions.
\end{ccnote}

\subsection{Text}

For a numeric cell this displays exactly as General does.

\section{Rounding on display}

A format rounds what is shown, not what is stored, and it rounds \textbf{half
away from zero}: at two places, 8.675 shows as \lit{8.68} and $-1.005$ shows as
\lit{-1.01}.

\begin{ccnote}
The underlying number is binary floating point, so a value that looks like
exactly 8.675 may be very slightly under it, and will then round down. If a
displayed total is one penny out from a hand-added column, this is why.
Appendix~\ref{app:numbers} explains the representation. \fn{ROUND} in a formula changes the
value itself, which is the fix when the value has to be exact.
\end{ccnote}

\section{Alignment}
\index{alignment}

Unless a cell carries an explicit alignment from a file, \ccprod{} decides for
itself:

\begin{cctable}{Cell holds}{Aligned}{2.2in}{2.62in}
A number & Right \\
A date & Right \\
Text & Left \\
\msg{TRUE} or \msg{FALSE} & Left \\
An error value & Left \\
A formula & By whatever the formula produced \\
\end{cctable}

\section{Column widths}
\index{columns!width}

The default width is ten characters. A width may be between three and forty,
and is stored per column per worksheet, but --- as with formats --- there is no
command that sets one. Widths arrive from \fname{.X16S} and \fname{.XLSX} files
and are saved back out again.

A value too wide for its column is cut off at the column edge, with no
\lit{\#\#\#\#} marker and no spill into the next cell.

\section{Dates}
\index{dates}

A date in \ccprod{} is a number: the count of days since the beginning of 1900,
with 1 being 1~January~1900. The date format is what makes it look like a date.

\subsection{The display}

One format only: \textbf{\lit{YYYY-MM-DD}}. There is no day-first or month-first
option, no month names, and no locale setting. The reason is that ISO order
sorts correctly as text and cannot be misread.

\subsection{Arithmetic}

Because a date is a number, arithmetic on it works: a date cell plus 30 is
thirty days later, and one date minus another is the number of days between
them.

\begin{cccaution}
The \emph{result} of that arithmetic is an ordinary cell in General format, so
it displays as a day number rather than as a date. \fml{=A1+30} on a date in
March 2023 shows \lit{45030}, not \lit{2023-04-14}. There is no way to put a
date format on the answer.
\end{cccaution}

\subsection{The 1900 leap year}

\begin{ccnote}
1900 was not a leap year. Excel believes it was, and every spreadsheet file
ever written carries the consequence: day number 60 is a 29~February that never
happened, and every date after it is one greater than the true count of days.

\ccprod{} \textbf{reproduces the error deliberately}, because reproducing it is
the only way to agree with the files it has to read. Day 60 displays as
\lit{1900-02-29}. Dates after 1~March~1900 are correct in every practical
sense; the discrepancy is entirely in the first two months of 1900.
\end{ccnote}

Day number 0 displays as \lit{1900-01-00}, which is not a date at all. It is
what the format produces for a value that predates the epoch, and it is a sign
that a cell has been formatted as a date by accident.

\subsection{Typing a date}

You cannot. There is no \fn{DATE} function, no \fn{TODAY}, and nothing that
parses \lit{2026-08-08} as anything but text. Date cells arrive only from
files. This is the largest single gap in the program's handling of dates and it
is worth planning around: if a worksheet needs dates in it, put them in the
file it is imported from.
